\begin{abstract}
Large language models (LLMs) are increasingly used as scoring and
representation engines for recommendation, yet comparisons across LLM-based
recommenders are often confounded by different candidate sets, backbones,
importers, and evaluation schemas. We first contribute a rigorous unified
same-candidate reranking protocol: every method ranks the same 101 candidates
(one positive and 100 popularity-sampled negatives) for each of 10{,}000 users,
uses the same Qwen3-8B backbone when an LLM is required, and is evaluated with a
shared importer, identical metrics, exact coverage checks, and paired
Holm-corrected bootstrap tests, with full provenance. Against eight
official-code-level LLM-based recommendation baselines across eight Amazon
domains (Beauty, Books, Electronics, Movies, Sports, Toys, Home, Tools), our
central empirical finding is that a \emph{task-grounded pointwise LLM relevance
posterior} -- the model's own per-candidate relevance probability, used directly
as the ranking score -- ranks first in six of the eight domains, improving
NDCG@10 by $+21.6\%$ to $+53.2\%$ over the strongest baseline (and on the six
winning domains it leads on essentially all seven metrics HR@5/@10/@20,
NDCG@5/@10/@20, MRR). In the remaining two domains it is competitive but not
first: rank~2 on Beauty ($-11\%$ NDCG@10 behind ProEx) and rank~5 on Movies
($-24\%$ behind LLMEmb). We report these two losses explicitly rather than
dropping them. On the four domains with full per-event signal rows, all 56
per-domain paired tests are positive and Holm-significant.

We then report a rigorous \emph{negative result} on the design space this work
set out to study. We instrument the pointwise posterior with a principled
calibrated-uncertainty and risk-adjusted ranking family that decomposes
uncertainty into boundary ambiguity, a calibration gap, evidence support, and
counterevidence. Under a leave-one-component-out ablation and a
raw-probability attribution, this machinery does not improve same-candidate
ranking: the zero-risk setting is test-best in all four diagnostic domains, a
confidence-only variant matches or exceeds the full family in three of those
four domains, and individual uncertainty terms are inert or mildly harmful. The
working ranking signal lies in the posterior itself, not in the
uncertainty decomposition. As a minor, descriptive observation -- and one we
caveat as partly circular -- the event-level uncertainty signal does stratify
reliability, but we do not claim it improves ranking. This is a scoped result
about same-candidate reranking, not a full-catalog recommender SOTA claim.
\end{abstract}
